% ---
% titre: Problème de proportionnalité - la provision de sucre
% theme: FA
% chapitre: Proportionnalité
% sous_chapitre: Résoudre des problèmes
% niveau: [10e]
% type: EVAL
% tags: [proportionnalite, p-question-TS]
% difficulte: 1
% corrige: true
% ---

\question[2]
Une réserve de sucre suffit à l'alimentation de 720 personnes pendant 24 jours. Que durerait la même provision s'il fallait nourrir 880 personnes ?

\vspace{10pt}{\footnotesize\raggedleft Espace pour ta démarche et tes calculs\par}
\begin{solutionorgrid}[140pt]

\vspace{5pt}
\def\arraystretch{1.5}
\begin{tabularx}{0.5\textwidth}{l|C|C|}
Temps [jours]&24&\\
\hline
Nombre de personnes&720&880\\
\end{tabularx}
\def\arraystretch{1}

\hfill\textbf{(0,5pt)} pour la construction du tableau

\vspace{15pt}$880 \cdot 24 : 720 = 29,$ francs \hfill calcul \textbf{(0,5pt)} et réponse \textbf{(0,5pt)}
\end{solutionorgrid}

\begin{tcolorbox}[enhanced jigsaw, interior hidden, colframe=box, parbox=false]%
Réponse: 
\sol{\vspace{20pt}}{\textbf{(0.5pt)}}
\end{tcolorbox}



